
\documentclass{article}
\usepackage{colortbl}
\usepackage{makecell}
\usepackage{multirow}
\usepackage{supertabular}

\begin{document}

\newcounter{utterance}

\centering \large Interaction Transcript for game `cladder', experiment `full\_v1.5\_default', episode 525 with Qwen3.5{-}9B.
\vspace{24pt}

{ \footnotesize  \setcounter{utterance}{1}
\setlength{\tabcolsep}{0pt}
\begin{supertabular}{c@{$\;$}|p{.15\linewidth}@{}p{.15\linewidth}p{.15\linewidth}p{.15\linewidth}p{.15\linewidth}p{.15\linewidth}}
    \# & \multicolumn{2}{c}{Player} && \multicolumn{2}{c}{Game Master} \\
    \hline

    \theutterance \stepcounter{utterance}  
    & & & \multicolumn{4}{p{0.6\linewidth}}{
        \cellcolor[rgb]{0.9,0.9,0.9}{
            \makecell[{{p{\linewidth}}}]{
                \texttt{\tiny{[P1$\langle$GM]}}
                \texttt{You are an expert in causal inference. The following question is not a typical commonsense query, but rather a meticulously designed question created by a professor specializing in causal inference, intended to assess the students' mastery of the course content.} \\
\texttt{Imagine a self{-}contained, hypothetical world with only the following conditions, and without any unmentioned factors or causal relationships: Unobserved confounders has a direct effect on drug taken and cholesterol level. Treatment assignment has a direct effect on drug taken. Drug taken has a direct effect on cholesterol level. Unobserved confounders is unobserved. For patients not assigned the drug treatment, the probability of low cholesterol is 63\%. For patients assigned the drug treatment, the probability of low cholesterol is 52\%. For patients not assigned the drug treatment, the probability of taking of all assigned drugs is 37\%. For patients assigned the drug treatment, the probability of taking of all assigned drugs is 74\%. Will taking of all assigned drugs increase the chance of low cholesterol?} \\
\texttt{Start your answer with 'Yes' or 'No', followed by additional reasoning or evidence to support your explanation.} \\
            }
        }
    }
     \\ \\

    \theutterance \stepcounter{utterance}  
    & \multicolumn{4}{p{0.6\linewidth}}{
        \cellcolor[rgb]{0.561,0.737,0.561}{
            \makecell[{{p{\linewidth}}}]{
                \texttt{\tiny{[P1$\rangle$GM]}}
                \texttt{No} \\
\\ 
\texttt{**Reasoning:**} \\
\\ 
\texttt{To determine if taking the drug increases the chance of low cholesterol, we must estimate the **Causal Effect of Treatment on the Treated (CETT)** or the **Average Treatment Effect (ATE)**, specifically isolating the effect of the *drug taken* from the *treatment assignment*.} \\
\\ 
\texttt{In this hypothetical world, the causal structure is defined as follows:} \\
\texttt{1.  \$U\$ (Unobserved Confounders) \$\textbackslash{}rightarrow\$ \$D\$ (Drug Taken)} \\
\texttt{2.  \$U\$ (Unobserved Confounders) \$\textbackslash{}rightarrow\$ \$Y\$ (Cholesterol Level)} \\
\texttt{3.  \$A\$ (Treatment Assignment) \$\textbackslash{}rightarrow\$ \$D\$ (Drug Taken)} \\
\texttt{4.  \$D\$ (Drug Taken) \$\textbackslash{}rightarrow\$ \$Y\$ (Cholesterol Level)} \\
\\ 
\texttt{The critical issue here is the presence of **unobserved confounding** (\$U\$). The variable \$U\$ affects both the decision to take the drug (\$D\$) and the outcome (cholesterol, \$Y\$). Because \$U\$ is unobserved, we cannot condition on it to block the backdoor path between \$D\$ and \$Y\$. Consequently, the observed association between drug taking and cholesterol is biased.} \\
\\ 
\texttt{Let's analyze the provided probabilities:} \\
\texttt{*   **Observation 1:** Among those *assigned* the drug (\$A=1\$), the probability of low cholesterol is 52\%.} \\
\texttt{*   **Observation 2:** Among those *not assigned* the drug (\$A=0\$), the probability of low cholesterol is 63\%.} \\
\\ 
\texttt{At first glance, one might be tempted to conclude that the drug reduces cholesterol (52\% vs 63\%). However, this comparison is flawed because the groups differ not only in treatment assignment but also in their propensity to take the drug and their underlying unobserved characteristics (\$U\$).} \\
\\ 
\texttt{*   **Group \$A=1\$ (Assigned):** 74\% take the drug.} \\
\texttt{*   **Group \$A=0\$ (Not Assigned):** 37\% take the drug.} \\
\\ 
\texttt{The group assigned to treatment is much more likely to take the drug. If the unobserved confounder \$U\$ is positively correlated with taking the drug (e.g., healthier or more motivated individuals are both more likely to be assigned or choose to take the drug, or perhaps the assignment mechanism correlates with \$U\$), then the group taking the drug might inherently have different baseline cholesterol levels due to \$U\$.} \\
\\ 
\texttt{Specifically, since \$U\$ affects both \$D\$ and \$Y\$, the correlation between \$D\$ and \$Y\$ observed in the data is a mixture of the causal effect of \$D\$ on \$Y\$ and the confounding effect of \$U\$. Without knowing the distribution of \$U\$ or the specific strength of the paths \$U \textbackslash{}rightarrow D\$ and} \\
            }
        }
    }
    & & \\ \\

    \theutterance \stepcounter{utterance}  
    & & & \multicolumn{2}{p{0.3\linewidth}}{
        \cellcolor[rgb]{0.9,0.9,0.9}{
            \makecell[{{p{\linewidth}}}]{
                \texttt{\tiny{[GM$|$GM]}}
                \texttt{No} \\
            }
        }
    }
    & & \\ \\

    \theutterance \stepcounter{utterance}  
    & & & \multicolumn{2}{p{0.3\linewidth}}{
        \cellcolor[rgb]{0.9,0.9,0.9}{
            \makecell[{{p{\linewidth}}}]{
                \texttt{\tiny{[GM$|$GM]}}
                \texttt{no} \\
            }
        }
    }
    & & \\ \\

    \theutterance \stepcounter{utterance}  
    & & & \multicolumn{2}{p{0.3\linewidth}}{
        \cellcolor[rgb]{0.9,0.9,0.9}{
            \makecell[{{p{\linewidth}}}]{
                \texttt{\tiny{[GM$|$GM]}}
                \texttt{game\_result = WIN} \\
            }
        }
    }
    & & \\ \\

\end{supertabular}
}

\end{document}
